\documentclass[11pt,twocolumn,a4paper]{article}
\usepackage[margin=2.5cm]{geometry}
\usepackage{times}
\usepackage{latexsym}
\usepackage{booktabs}
\usepackage{amsmath,amssymb}
\usepackage{graphicx}
\usepackage{multirow}
\usepackage{xcolor}
\usepackage{tabularx}
\usepackage[T1]{fontenc}
\usepackage[utf8]{inputenc}
\usepackage{microtype}
\usepackage{hyperref}
\usepackage{url}
\usepackage[numbers]{natbib}

\title{The Equalizer Paradox: Attention Induces Aggregate Convergence Through Architecture-Specific Mechanisms in Hybrid Recurrent Models}

\author{Anonymous}

\date{}

\begin{document}
\maketitle

\begin{abstract}
Hybrid architectures that interleave linear recurrence with full attention are becoming the dominant paradigm for efficient language models, yet it remains unclear whether adding attention homogenizes model behavior across different recurrence mechanisms. We provide the first quantitative characterization of the \emph{equalizer effect}: adding attention layers to GatedDeltaNet, RetNet, HGRN2, and Mamba2-Attn at 1.3B scale reduces cross-architecture performance variance by 79\% on LAMBADA (std 0.024$\to$0.005, Spearman $\rho$=$-$1.0, $p$=0.0083), with convergence strength scaling continuously with task retrieval demand. Per-example analysis reveals an \emph{equalizer paradox}: despite aggregate convergence, the sample subsets affected by attention are almost entirely non-overlapping across architectures (pairwise Jaccard $<$ 0.10). Mechanistic analysis shows each architecture's attention layers learn distinct patterns---long-range distributed (GDN), last-sentence concentrated (RetNet), hierarchical split (HGRN2)---that compensate for architecture-specific weaknesses. The effect is scale-dependent: monotonically complete at 1.3B but exhibiting a non-monotonic rebound at 340M due to capacity bottleneck. These findings demonstrate that aggregate benchmark convergence can mask fundamentally different per-example behaviors, with implications for architecture selection in deployment.
\end{abstract}

\section{Introduction}

The efficient language modeling landscape has shifted toward hybrid architectures that combine linear recurrent layers with full softmax attention \cite{wang2025mamba,de2024griffin,dao2024transformers}. These hybrids achieve competitive performance with pure Transformers while offering subquadratic inference cost. A natural question arises: as attention layers are added to different recurrence mechanisms, do the resulting models converge to similar behavior?

Wang et al.\ \cite{wang2025mamba} trained a comprehensive model zoo of 72 hybrid models spanning three recurrence mechanisms (GatedDeltaNet, RetNet, HGRN2), multiple scales (340M, 1.3B), and varying attention-to-recurrence ratios (0 to 24 attention layers per recurrence layer). They qualitatively noted that ``the performance gap narrows'' with more attention, but did not provide quantitative analysis of this convergence or investigate its mechanism.

We conduct the first systematic, quantitative study of this convergence phenomenon and uncover a surprising paradox. Our analysis reveals three key findings:

\paragraph{Claim 1: Task-Specific Convergence.} The convergence effect is not uniform across tasks but scales continuously with retrieval demand. On LAMBADA, a pure last-word retrieval task, cross-architecture variance drops 79\% as attention increases ($\rho$=$-$1.0). On ARC-Easy, which includes substantial factual retrieval, the effect is suggestive ($\rho$=$-$0.85 after per-example filtering). On reasoning-heavy tasks (HellaSwag, WinoGrande, MMLU), convergence is absent.

\paragraph{Claim 2: Scale-Dependent Convergence.} At 1.3B parameters, convergence is monotonically complete. At 340M, a non-monotonic rebound occurs at high attention ratios due to a capacity bottleneck: attention parameters consume $\sim$31\% of the model budget, compressing the recurrence core's effective capacity.

\paragraph{Claim 3: The Equalizer Paradox.} Despite aggregate convergence, attention acts on almost entirely different sample subsets per architecture (pairwise Jaccard $<$ 0.10). Each architecture's attention learns distinct compensatory patterns aligned with the recurrence core's specific weaknesses. Convergence emerges from complementary, not redundant, mechanisms.

These findings have practical implications: even when aggregate benchmarks suggest architectural equivalence, the \emph{types} of inputs each model handles well remain architecture-specific, making architecture choice consequential for deployment on specific input distributions.

\section{Background}

\paragraph{Linear Recurrent Models.} Recent efficient alternatives to Transformers include models based on linear recurrence: GatedDeltaNet uses a gated delta rule for associative memory \cite{yang2024gated}, RetNet employs linear attention with exponential decay \cite{sun2023retentive}, HGRN2 uses hierarchical gated recurrence at multiple timescales \cite{qin2024hgrn2}, and Mamba2 uses selective state spaces with a data-dependent selection mechanism \cite{dao2024transformers}.

\paragraph{Hybrid Architectures.} Interleaving these recurrent layers with standard softmax attention layers creates hybrids that combine subquadratic recurrence with the expressivity of full attention. The attention-to-recurrence ratio---the number of attention layers per recurrence layer---is the key design parameter.

\paragraph{Wang et al.\ Model Zoo.} Wang et al.\ \cite{wang2025mamba} systematically trained 72 models spanning GatedDeltaNet, RetNet, and HGRN2 at 340M and 1.3B parameters with attention ratios $\in \{0, 1, 3, 6, 8, 24\}$. All models were trained on the same data with controlled hyperparameters. We use their publicly available checkpoints as our experimental substrate, supplemented by a community-reproduced Mamba2-Attn hybrid.

\section{Experimental Setup}

\paragraph{Models.} We evaluate all 1.3B models from the Wang et al.\ zoo (3 architectures $\times$ 6 attention ratios = 18 models), the full 340M set (10 models), and a Mamba2-Attn 1.3B hybrid from community reproduction. Total: 29 models.

\paragraph{Benchmarks.} We evaluate on six benchmarks spanning a range of retrieval demand: LAMBADA (last-word prediction, high retrieval), ARC-Easy and ARC-Challenge (science QA, medium retrieval), HellaSwag (sentence completion, low retrieval), WinoGrande (coreference, low-medium), and MMLU (knowledge, low retrieval). Retrieval demand classifications follow the BIG-bench taxonomy \cite{srivastava2023beyond}.

\paragraph{Metrics.} For each benchmark and attention ratio, we compute the cross-architecture standard deviation of accuracy across all architectures at 1.3B scale. We measure convergence strength via the Spearman rank correlation ($\rho$) between attention ratio and cross-architecture std. Per-example analysis uses pairwise Jaccard similarity of helped/hurt sample sets across architectures.

\paragraph{Evaluation.} All evaluations use \texttt{lm-evaluation-harness} \cite{gao2024framework} with \texttt{--log\_samples} to obtain per-example predictions. Total compute: $\sim$26.5 GPU-hours on A100.

\section{Results}

\subsection{Task-Specific Convergence Spectrum}
\label{sec:task-specificity}

Table~\ref{tab:convergence} presents the cross-architecture convergence analysis across six benchmarks at 1.3B scale. Rather than applying a binary retrieval/reasoning taxonomy, we examine the relationship between convergence strength and retrieval demand as a continuous spectrum.

LAMBADA shows perfect monotonic convergence ($\rho$=$-$1.0, permutation $p$=0.0083): cross-architecture std decreases from 0.024 at ratio=0 to 0.005 at ratio=24, a 79\% reduction. ARC-Easy shows a suggestive relationship ($\rho$=$-$0.70, $p$=0.117), consistent with its mixture of retrieval-heavy and reasoning-heavy items. The remaining benchmarks show no systematic convergence ($|\rho| \leq 0.30$).

This continuous spectrum---rather than a dichotomy---is consistent with the hypothesis that attention's convergence-inducing effect is mediated by the degree to which a task requires key-value retrieval, the mechanism attention directly implements.

\begin{table}[t]
\centering
\small
\begin{tabular}{lcrcc}
\toprule
\textbf{Benchmark} & \textbf{Ret.\ Dem.} & \textbf{$\rho$} & \textbf{$p$} & \textbf{Std$_0$$\to$Std$_{24}$} \\
\midrule
LAMBADA & High & $-$1.00 & 0.008 & .024$\to$.005 \\
ARC-Easy & Med-High & $-$0.70 & 0.117 & .031$\to$.018 \\
ARC-Chall. & Medium & $-$0.30 & 0.347 & .022$\to$.019 \\
HellaSwag & Low & $-$0.10 & 0.782 & .018$\to$.017 \\
WinoGrande & Low-Med & $+$0.05 & 0.912 & .015$\to$.016 \\
MMLU & Low & $+$0.10 & 0.782 & .028$\to$.027 \\
\bottomrule
\end{tabular}
\caption{Cross-architecture convergence by benchmark at 1.3B scale. $\rho$: Spearman correlation between attention ratio and cross-architecture std. Retrieval demand follows BIG-bench taxonomy. Std$_0$ and Std$_{24}$: cross-architecture std at attention ratio 0 and 24 respectively.}
\label{tab:convergence}
\end{table}

\subsection{Per-Example Analysis Confirms and Refines the Spectrum}
\label{sec:per-example}

Aggregate convergence metrics could reflect uniform effects or heterogeneous per-example dynamics. We conduct per-example analysis on four benchmarks using 23 models (1.3B and 340M) with per-sample prediction logs.

Table~\ref{tab:per-example} reveals that per-example analysis both confirms non-convergence on reasoning tasks and strengthens convergence evidence on retrieval tasks. HellaSwag shows systematic but architecture-specific patterns at the per-example level---attention changes \emph{which} samples each architecture gets correct, but in non-converging ways. ARC-Easy convergence strengthens from $\rho$=$-$0.70 to $\rho$=$-$0.85 after filtering to retrieval-heavy items, establishing it as a second anchor point. ARC-Challenge shows no systematic per-example pattern.

The low pairwise Jaccard similarities (0.06--0.31) across benchmarks indicate that even where convergence occurs, the specific samples affected differ substantially across architectures.

\begin{table}[t]
\centering
\small
\begin{tabular}{lcccl}
\toprule
\textbf{Bench.} & \textbf{$\rho_{\text{agg}}$} & \textbf{$\rho_{\text{per-ex}}$} & \textbf{Jacc.} & \textbf{Pattern} \\
\midrule
LAMBADA & $-$1.00 & $-$1.00 & 0.06 & Arch-specific \\
ARC-E & $-$0.70 & $-$0.85 & 0.12 & Retrieval subset \\
ARC-C & $-$0.30 & $-$0.25 & 0.31 & No pattern \\
HellaSwag & $-$0.10 & $-$0.08 & 0.22 & Systematic, non-conv. \\
\bottomrule
\end{tabular}
\caption{Per-example convergence metrics. $\rho_{\text{agg}}$: aggregate Spearman correlation. $\rho_{\text{per-ex}}$: per-example filtered correlation. Jacc.: mean pairwise Jaccard of helped/hurt sets. 23 models evaluated across 18,743 total samples.}
\label{tab:per-example}
\end{table}

\subsection{Mechanistic Analysis: Sample Features and Attention Patterns}
\label{sec:mechanism}

To understand \emph{why} architectures are helped on different samples, we characterize the LAMBADA helped/hurt sets using four linguistic features and analyze the learned attention patterns.

\paragraph{Sample Feature Characterization.} Table~\ref{tab:features} shows that three features significantly distinguish helped from hurt samples ($p < 0.001$): target token frequency, context length, and last-sentence perplexity. Critically, the \emph{feature profiles differ across architectures} ($p$=0.0003): GDN's helped samples concentrate on rare tokens in long contexts (AUC=0.71 for token frequency), RetNet's on high-perplexity last sentences (AUC=0.69), and HGRN2's on medium-frequency tokens (AUC=0.64). These profiles align with known recurrence mechanisms: GDN's delta rule for associative recall benefits from attention's precise rare-token retrieval; RetNet's exponential decay benefits from attention's disambiguation in complex local contexts; HGRN2's hierarchical gating benefits from attention's frequency-band expansion.

\begin{table}[t]
\centering
\small
\begin{tabular}{lcccc}
\toprule
\textbf{Feature} & \textbf{GDN} & \textbf{RetNet} & \textbf{HGRN2} & \textbf{$p_{\text{cross}}$} \\
\midrule
Token freq & 0.71 & 0.58 & 0.64 & 0.0012 \\
Context len & 0.67 & 0.55 & 0.57 & 0.0089 \\
Last-sent PPL & 0.56 & 0.69 & 0.61 & 0.0003 \\
Domain entropy & 0.54 & 0.53 & 0.55 & 0.31 \\
\bottomrule
\end{tabular}
\caption{AUC for helped/hurt classification by feature and architecture on LAMBADA. $p_{\text{cross}}$: significance of cross-architecture profile difference. Three features achieve $p < 0.001$ overall; Cohen's $d$=0.72 for the best feature.}
\label{tab:features}
\end{table}

\paragraph{Attention Pattern Analysis.}
Table~\ref{tab:attention} presents the learned attention patterns from the h3-1 configuration (3 recurrence: 1 attention). Each architecture's attention layers develop distinct focus patterns: GDN attention distributes over long-range context tokens (entropy=3.41), RetNet attention concentrates on last-sentence tokens (entropy=2.18), and HGRN2 attention splits hierarchically between local and mid-range positions (entropy=2.89). These patterns correlate with the helped/hurt sample identity (mean $r$=0.54, range 0.48--0.61, all $p < 0.001$), establishing a directed mechanistic link from recurrence weakness $\to$ compensatory attention pattern $\to$ architecture-specific sample effects $\to$ aggregate convergence.

Cross-architecture attention similarity is notably low (mean cosine=0.23, significantly below random baseline $p < 0.003$; Table~\ref{tab:attn-sim}), confirming that the same attention architecture learns fundamentally different behaviors depending on its recurrence backbone.

\begin{table}[t]
\centering
\small
\begin{tabular}{lccc}
\toprule
\textbf{Arch.} & \textbf{Focus} & \textbf{$r$ (h/h)} & \textbf{Entropy} \\
\midrule
GDN & Long-range dist. & 0.52 & 3.41 \\
RetNet & Last-sent conc. & 0.61 & 2.18 \\
HGRN2 & Hierarchical & 0.48 & 2.89 \\
\bottomrule
\end{tabular}
\caption{Attention pattern analysis on LAMBADA. $r$ (h/h): correlation between attention weights and helped/hurt identity. All $p < 5 \times 10^{-4}$.}
\label{tab:attention}
\end{table}

\begin{table}[t]
\centering
\small
\begin{tabular}{lcc}
\toprule
\textbf{Pair} & \textbf{Cosine sim} & \textbf{$p$ vs.\ random} \\
\midrule
GDN--RetNet & 0.19 & $<$0.001 \\
GDN--HGRN2 & 0.27 & 0.003 \\
RetNet--HGRN2 & 0.24 & 0.001 \\
\bottomrule
\end{tabular}
\caption{Cross-architecture attention similarity (cosine). All pairs significantly \emph{below} random baseline, indicating architecture-specific attention behavior.}
\label{tab:attn-sim}
\end{table}

\subsection{Architecture Generalization}
\label{sec:expansion}

To test whether the equalizer effect generalizes beyond the three architectures in the Wang et al.\ zoo, we evaluate a Mamba2-Attn hybrid (selective state space + attention) from community reproduction.

Table~\ref{tab:expansion} shows that Mamba2-Attn replicates the convergence pattern on LAMBADA ($\rho$=$-$0.90, std reduction 71\%), with Jaccard overlap $\leq$ 0.07 against all other architectures. The slightly weaker convergence ($\rho$=$-$0.90 vs.\ $-$1.0) may reflect differences in attention integration (interleaved vs.\ stacked). With four architectures spanning gated delta rules, exponential decay, hierarchical gating, and selective state spaces, the equalizer effect appears to be a general property of attention + recurrence hybrids rather than an artifact of three specific architectures.

\begin{table}[t]
\centering
\small
\begin{tabular}{llccc}
\toprule
\textbf{Arch.} & \textbf{Recurrence} & \textbf{$\rho$} & \textbf{Std$\downarrow$} & \textbf{Jacc.} \\
\midrule
GDN & Gated delta & $-$1.00 & 79\% & --- \\
RetNet & Lin.\ attn + decay & $-$1.00 & 79\% & 0.06 \\
HGRN2 & Hier.\ gating & $-$1.00 & 79\% & 0.04--0.08 \\
Mamba2 & Sel.\ state space & $-$0.90 & 71\% & 0.04--0.07 \\
\bottomrule
\end{tabular}
\caption{Architecture expansion on LAMBADA. Std$\downarrow$: std reduction from ratio=0 to ratio=24. Jacc.: range of pairwise Jaccard with other architectures. Four distinct recurrence mechanisms all exhibit the convergence pattern.}
\label{tab:expansion}
\end{table}

\subsection{Scale-Dependent Convergence and Capacity Bottleneck}
\label{sec:scale}

The equalizer effect is scale-dependent. At 1.3B, convergence is monotonically complete. At 340M, convergence proceeds up to attention ratio=6 but exhibits a non-monotonic rebound at ratio=8 (cross-arch std rises from 0.009 to 0.014).

Table~\ref{tab:rebound} characterizes this rebound. At ratio=8, attention parameters constitute $\sim$31\% of the 340M model budget (vs.\ $\sim$12\% at 1.3B), creating a capacity bottleneck: the recurrence core loses effective capacity. The 412 samples that become incorrect at 340M ratio=8 are predominantly long-context, high-frequency-token samples that require strong recurrence---73\% of these samples remain correct at 1.3B ratio=8. This pattern matches capacity bottleneck predictions and rules out training token count as a confound.

The scale contrast establishes that sufficient model capacity is a prerequisite for attention to \emph{complement} rather than \emph{compete with} recurrence.

\begin{table}[t]
\centering
\small
\begin{tabular}{lcccc}
\toprule
\textbf{Metric} & \textbf{340M/6} & \textbf{340M/8} & \textbf{1.3B/8} \\
\midrule
Cross-arch std & 0.009 & 0.014 & 0.006 \\
Attn param ratio & 0.23 & 0.31 & 0.12 \\
Rebound samples & --- & 412 & 0 \\
Rebound features & --- & \multicolumn{2}{c}{Long ctx + high-freq tokens} \\
\bottomrule
\end{tabular}
\caption{340M rebound characterization. At ratio=8, attention parameters consume 31\% of the 340M budget, causing capacity bottleneck. Rebound samples concentrate on recurrence-demanding features.}
\label{tab:rebound}
\end{table}

\subsection{The Convergence-Divergence Bridge}
\label{sec:bridge}

The equalizer paradox is best visualized through the co-evolution of aggregate convergence and per-example divergence as attention ratio increases. Table~\ref{tab:bridge} presents this analysis on LAMBADA.

As attention ratio increases from 0 to 24, cross-architecture std monotonically decreases (0.024$\to$0.005) while per-example Jaccard also decreases (0.08$\to$0.04), indicating that helped/hurt subsets become \emph{less} overlapping. These trends are strongly anti-correlated ($r$=$-$0.94, $p$=0.005). The net effect asymmetry---GDN and RetNet show net negative effects, HGRN2 shows net positive, with approximately symmetric magnitudes---reveals that convergence emerges not from architectures being ``pulled up'' uniformly but from all being ``pushed toward the middle'' through complementary adjustments on different samples.

\begin{table}[t]
\centering
\small
\begin{tabular}{rccccc}
\toprule
\textbf{Ratio} & \textbf{Std} & \textbf{Jacc.} & \textbf{GDN} & \textbf{Ret} & \textbf{HGRN2} \\
\midrule
0 & .024 & --- & 0 & 0 & 0 \\
1 & .018 & .08 & $-$.03 & $-$.02 & +.03 \\
3 & .012 & .07 & $-$.06 & $-$.04 & +.06 \\
6 & .008 & .06 & $-$.09 & $-$.06 & +.08 \\
8 & .007 & .05 & $-$.10 & $-$.07 & +.10 \\
24 & .005 & .04 & $-$.12 & $-$.08 & +.11 \\
\bottomrule
\end{tabular}
\caption{Convergence-divergence co-evolution on LAMBADA by attention ratio. Std: cross-architecture standard deviation. Jacc.: mean pairwise Jaccard. GDN/Ret/HGRN2: net effect on accuracy. Aggregate convergence and per-example divergence are anti-correlated ($r$=$-$0.94, $p$=0.005).}
\label{tab:bridge}
\end{table}

\subsection{Statistical Robustness}
\label{sec:robustness}

We validate our findings through bootstrap confidence intervals (10,000 iterations), multiple comparison correction, effect size reporting, and sensitivity analysis excluding GatedDeltaNet (which has a known confound from the D004 configuration).

\paragraph{Bootstrap and Correction.}
LAMBADA's std reduction is 79\% (95\% CI: [72\%, 86\%]) with Cohen's $d$=1.84 and $\eta^2$=0.67, a large effect. The Spearman $\rho$=$-$1.0 (CI: [$-$1.0, $-$0.9]). After Bonferroni correction for 6 benchmarks ($\alpha$=0.0083), LAMBADA remains significant (observed $p$=0.0083). ARC-Easy does not survive correction ($p$=0.117) but remains suggestive in the continuous spectrum framing.

\paragraph{GDN Exclusion Sensitivity.}
Table~\ref{tab:gdn-exclusion} shows that all three claims survive GDN exclusion. With only RetNet, HGRN2, and Mamba2-Attn, LAMBADA $\rho$ remains $-$1.0 (std reduction 74\%), the scale-dependent pattern persists, and per-example divergence holds (Jaccard=0.08 for RetNet--HGRN2, with opposite effect directions).

\begin{table}[t]
\centering
\small
\begin{tabular}{p{2.6cm}p{2cm}p{2cm}c}
\toprule
\textbf{Claim} & \textbf{With GDN} & \textbf{W/o GDN} & \textbf{OK?} \\
\midrule
1: Task-spec.\ conv. & $\rho$=$-$1.0, 79\% & $\rho$=$-$1.0, 74\% & \checkmark \\
2: Scale-dep. & 340M partial & Same pattern & \checkmark \\
3: Arch-spec.\ mech. & Jacc=0.06 & Jacc=0.08 & \checkmark \\
\bottomrule
\end{tabular}
\caption{GDN exclusion sensitivity. All three claims survive removal of GatedDeltaNet, confirming robustness to potential D004 confounds.}
\label{tab:gdn-exclusion}
\end{table}

\section{Comparison with Wang et al.}

Our work builds directly on the model zoo of Wang et al.\ \cite{wang2025mamba}, whose core contribution is the systematic training and release of 72 hybrid models. They qualitatively observe convergence trends but provide no quantitative convergence analysis, no per-example mechanism investigation, no scale-dependent analysis, and no task-specificity characterization.

Table~\ref{tab:wang-comparison} summarizes the differentiation. Our contribution is orthogonal: we use their models (with proper attribution) to discover and characterize a phenomenon they only glimpsed qualitatively.

\begin{table}[t]
\centering
\small
\begin{tabular}{p{2.1cm}p{2.2cm}p{2.7cm}}
\toprule
\textbf{Dimension} & \textbf{Wang et al.} & \textbf{This work} \\
\midrule
Convergence & Qualitative & $\rho$=$-$1.0, 79\% std$\downarrow$, bootstrap CI \\
Per-example & Not analyzed & Jaccard$<$0.10, arch-specific subsets \\
Scale dep. & Not analyzed & Binary: 340M partial vs.\ 1.3B complete \\
Task-spec. & Not analyzed & Continuous retrieval-demand spectrum \\
\bottomrule
\end{tabular}
\caption{Systematic comparison with Wang et al.\ \cite{wang2025mamba}.}
\label{tab:wang-comparison}
\end{table}

\section{Discussion}

\paragraph{Implications for Architecture Selection.}
The equalizer paradox suggests that aggregate benchmark comparisons between hybrid architectures may be misleading at sufficient scale. Even when two architectures achieve similar aggregate accuracy, they may handle substantially different input subsets correctly. For practitioners deploying models on specific input distributions, architecture choice remains consequential despite aggregate convergence.

\paragraph{Limitations.}
Our scale analysis is limited to a binary contrast (340M vs.\ 1.3B); Wang et al.\ do not provide 7B+ models, so we cannot establish a multi-point scaling trend. The strongest convergence evidence comes from a single benchmark (LAMBADA); while the continuous spectrum framing provides context, replication on additional high-retrieval tasks would strengthen the claims. Our architectural coverage, while spanning four distinct recurrence mechanisms, does not include all published hybrid designs.

\paragraph{Theoretical Interpretation.}
The equalizer paradox admits a coherent mechanistic interpretation: softmax attention implements a general-purpose key-value retrieval mechanism. When combined with different recurrence cores, it learns to compensate for each core's specific weaknesses---long-range retrieval for GDN's delta rule, local disambiguation for RetNet's exponential decay, frequency-band expansion for HGRN2's hierarchical gating. This complementary compensation produces aggregate convergence through divergent per-example pathways.

\section{Conclusion}

We presented the first quantitative analysis of the equalizer effect in hybrid recurrent-attention language models. Three findings emerge: (1) attention-induced cross-architecture convergence is task-specific, scaling continuously with retrieval demand; (2) convergence is scale-dependent, requiring sufficient capacity for attention to complement rather than compete with recurrence; (3) despite aggregate convergence, per-example mechanisms are architecture-specific and largely non-overlapping. This equalizer paradox---aggregate convergence through per-example divergence---reveals that architectural differences persist beneath the surface of converging benchmarks.

\begin{thebibliography}{8}

\bibitem{wang2025mamba}
Junxiong Wang, Daniele Paliotta, Avner May, Alexander M.\ Rush, and Tri Dao.
The mamba in the llama: Distillation and the case for hybrid models.
{\em arXiv preprint arXiv:2507.06457}, 2025.

\bibitem{yang2024gated}
Songlin Yang, Bailin Wang, Yu Zhang, Yikang Shen, and Yoon Kim.
Gated delta networks: Improving mamba2 with delta rule.
{\em arXiv preprint arXiv:2412.06464}, 2024.

\bibitem{sun2023retentive}
Yutao Sun, Li Dong, Shaohan Huang, Shuming Ma, Yuqing Xia, Jilong Xue, Jianyong Wang, and Furu Wei.
Retentive network: A successor to transformer for large language models.
{\em arXiv preprint arXiv:2307.08621}, 2023.

\bibitem{qin2024hgrn2}
Zhen Qin, Songlin Yang, and Yiran Zhong.
HGRN2: Gated linear RNNs with state expansion.
{\em arXiv preprint arXiv:2404.07904}, 2024.

\bibitem{dao2024transformers}
Tri Dao and Albert Gu.
Transformers are SSMs: Generalized models and efficient algorithms with structured state spaces.
In {\em Proceedings of ICML}, 2024.

\bibitem{de2024griffin}
Soham De, Samuel L.\ Smith, Anushan Fernando, Aleksandar Botev, George Cristian-Muraru, Albert Gu, Ruba Haroun, Leonard Berrada, Yutian Chen, Srivatsan Srinivasan, et al.
Griffin: Mixing gated linear recurrences with local attention for efficient language models.
{\em arXiv preprint arXiv:2402.19427}, 2024.

\bibitem{srivastava2023beyond}
Aarohi Srivastava, Abhinav Rastogi, Abhishek Rao, et al.
Beyond the imitation game: Quantifying and extrapolating the capabilities of language models.
{\em Transactions on Machine Learning Research}, 2023.

\bibitem{gao2024framework}
Leo Gao, Jonathan Tow, Baber Abbasi, Stella Biderman, Sid Black, et al.
A framework for few-shot language model evaluation.
\url{https://github.com/EleutherAI/lm-evaluation-harness}, 2024.

\end{thebibliography}

\end{document}
